% !TEX program = xelatex
% Lernplatt - by Can Siebert
% Kurs 71 (Dig 42) - Tag 5 - Loesungsheft

\documentclass[11pt,a4paper,oneside]{article}

\usepackage{polyglossia}
\setdefaultlanguage{german}
\usepackage[a4paper,margin=2.5cm]{geometry}

\usepackage{fontspec}
\defaultfontfeatures{Ligatures=TeX}
\IfFontExistsTF{Charter}{\setmainfont{Charter}}{%
  \IfFontExistsTF{Bitstream Charter}{\setmainfont{Bitstream Charter}}{\setmainfont{TeX Gyre Schola}}%
}
\IfFontExistsTF{Source Sans 3}{\setsansfont{Source Sans 3}}{%
  \IfFontExistsTF{Source Sans Pro}{\setsansfont{Source Sans Pro}}{\setsansfont{TeX Gyre Heros}}%
}
\newcommand{\caveatfontfallback}{\itshape}
\IfFontExistsTF{Caveat}{\newfontfamily\caveatfont{Caveat}[Scale=1.15]}{\newcommand\caveatfont{\caveatfontfallback}}
\setmonofont{Latin Modern Mono}

\usepackage{microtype}
\linespread{1.15}

\usepackage{xcolor}
\definecolor{lpInk}{RGB}{20,28,38}
\definecolor{lpAccent}{RGB}{157,74,26}
\definecolor{lpAccentLight}{RGB}{246,228,212}
\definecolor{lpAmber}{RGB}{222,138,30}
\definecolor{lpAmberLight}{RGB}{255,243,217}
\definecolor{lpAmberBorder}{RGB}{196,118,18}
\definecolor{lpGray}{RGB}{96,104,114}
\definecolor{lpGrayLight}{RGB}{238,240,243}
\definecolor{lpGreen}{RGB}{34,120,72}
\definecolor{lpGreenLight}{RGB}{222,238,228}
\definecolor{lpGreenBorder}{RGB}{24,96,58}
\definecolor{lpL1}{RGB}{34,120,72}
\definecolor{lpL2}{RGB}{22,90,140}
\definecolor{lpL3}{RGB}{170,42,52}

\usepackage[most]{tcolorbox}
\tcbuselibrary{breakable,skins}

\newtcolorbox{infobox}[1][Hinweis]{%
  enhanced, breakable, colback=lpAccentLight, colframe=lpAccent,
  boxrule=0.6pt, arc=2pt, left=10pt,right=10pt,top=8pt,bottom=8pt,
  fonttitle=\sffamily\bfseries\color{white}, coltitle=white,
  title={#1}, attach boxed title to top left={xshift=8pt,yshift=-8pt},
  boxed title style={size=small, colback=lpAccent, colframe=lpAccent, boxrule=0.4pt, arc=1pt},
  top=14pt
}

\newtcolorbox{loesungbox}[1][Musterl\"osung]{%
  enhanced, breakable, colback=lpGreenLight, colframe=lpGreenBorder,
  boxrule=0.6pt, arc=2pt, left=10pt,right=10pt,top=8pt,bottom=8pt,
  fonttitle=\sffamily\bfseries\color{white}, coltitle=white,
  title={#1}, attach boxed title to top left={xshift=8pt,yshift=-8pt},
  boxed title style={size=small, colback=lpGreenBorder, colframe=lpGreenBorder, boxrule=0.4pt, arc=1pt},
  top=14pt
}

\usepackage{enumitem}
\setlist{itemsep=2pt, topsep=4pt, parsep=0pt}
\setlist[itemize,1]{label={\textbullet},leftmargin=1.6em}
\setlist[itemize,2]{label={\textendash},leftmargin=1.4em}
\setlist[enumerate,1]{label=\arabic*., leftmargin=1.8em}

\usepackage{tabularx}
\usepackage{array}
\usepackage{booktabs}
\usepackage{longtable}

\usepackage{fancyhdr}
\usepackage{lastpage}
\pagestyle{fancy}
\fancyhf{}
\renewcommand{\headrulewidth}{0.3pt}
\renewcommand{\footrulewidth}{0pt}
\fancyhead[L]{\footnotesize\sffamily\color{lpGray}L\"osungsheft \textperiodcentered{} Kurs 71 \textperiodcentered{} Tag 5}
\fancyhead[R]{\footnotesize\sffamily\color{lpGray}Visio/Lucidchart \textperiodcentered{} EA \textperiodcentered{} MDA}
\fancyfoot[C]{\footnotesize\sffamily\color{lpGray}\thepage{} / \pageref{LastPage}}

\fancypagestyle{plain}{%
  \fancyhf{}%
  \renewcommand{\headrulewidth}{0pt}%
  \fancyfoot[C]{\footnotesize\sffamily\color{lpGray}\thepage{} / \pageref{LastPage}}%
}

\usepackage[explicit]{titlesec}
\titleformat{\section}[block]
  {\sffamily\Large\bfseries\color{lpAccent}}
  {\thesection.}{0.6em}{#1}
  [{\vspace{2pt}\color{lpAccent}\titlerule[0.6pt]}]
\titleformat{\subsection}[block]
  {\sffamily\large\bfseries\color{lpInk}}
  {\thesubsection}{0.6em}{#1}
\titlespacing*{\section}{0pt}{14pt}{8pt}
\titlespacing*{\subsection}{0pt}{10pt}{4pt}

\usepackage[hidelinks,unicode]{hyperref}

\newcommand{\akzent}[1]{{\caveatfont\color{lpAmber}#1}}
\newcommand{\fachbegriff}[1]{\textbf{#1}}
\newcommand{\quelle}[1]{{\footnotesize\sffamily\color{lpGray}(#1)}}

\newcommand{\niveaubadge}[2]{%
  \tcbox[on line, colback=#1, colframe=#1, coltext=white,
         boxrule=0pt, arc=2pt, left=5pt,right=5pt,top=1pt,bottom=1pt,
         nobeforeafter]{\sffamily\bfseries\footnotesize #2}%
}
\newcommand{\nivLi}{\niveaubadge{lpL1}{L1\,\textbullet\,Wissen}}
\newcommand{\nivLii}{\niveaubadge{lpL2}{L2\,\textbullet\,Anwendung}}
\newcommand{\nivLiii}{\niveaubadge{lpL3}{L3\,\textbullet\,Management}}

\newcommand{\zeit}[1]{%
  \tcbox[on line, colback=lpGrayLight, colframe=lpGrayLight, coltext=lpInk,
         boxrule=0pt, arc=2pt, left=5pt,right=5pt,top=1pt,bottom=1pt,
         nobeforeafter]{\sffamily\footnotesize\textbf{$\sim$\,#1}}%
}

\newcommand{\aufgabenkopf}[4]{%
  \par\vspace{6pt}
  \noindent{\sffamily\Large\bfseries\color{lpAccent}Aufgabe #1 \textendash{} #2}\par
  \vspace{2pt}
  \noindent{\color{lpAccent}\rule{\linewidth}{0.6pt}}\par
  \vspace{4pt}
  \noindent #3 \hfill \zeit{#4}\par
  \vspace{6pt}
}

\newcommand{\ytquery}[1]{%
  \par\vspace{4pt}
  \noindent{\footnotesize\sffamily\color{lpGray}\textbf{YouTube-Suchquery:} \texttt{#1}}\par
}

\begin{document}

\begin{titlepage}
  \pagestyle{empty}
  \vspace*{1.5cm}
  \noindent{\sffamily\footnotesize\color{lpGray}LERNPLATT \textperiodcentered{} BY CAN SIEBERT}\\[2pt]
  \noindent{\color{lpAccent}\rule{\linewidth}{1.2pt}}\\[4pt]
  \noindent{\sffamily\footnotesize\color{lpGray}L\"OSUNGSHEFT \textperiodcentered{} MODUL 2 \textperiodcentered{} TAG 5 VON 16 \textperiodcentered{} KURS 71 (Dig 42) \textperiodcentered{} 3.\,Juni 2026}

  \vspace{3cm}
  {\sffamily\fontsize{30}{34}\selectfont\bfseries\color{lpInk}L\"osungsheft\\Tag 5\par}

  \vspace{0.7cm}
  {\sffamily\large\color{lpGray}Musterl\"osungen \textendash{} Visio/Lucidchart,\\EA-Rahmen und MDA-Steuerungslogik.\par}

  \vspace{1.5cm}
  {\caveatfont\LARGE\color{lpGreen}Musterl\"osungen \textperiodcentered{} nicht die einzige Antwort}\par

  \vfill

  \noindent\begin{minipage}{0.55\linewidth}
    {\sffamily\footnotesize\color{lpGray}Hinweis:}\\
    {\sffamily\small\color{lpInk}Musterl\"osungen sind Diskussionsbasis,\\nicht die einzige korrekte L\"osung.}\\[10pt]
    {\sffamily\footnotesize\color{lpGray}Begleitend:}\\
    {\sffamily\small\color{lpInk}Aufgabenheft \& Lehrskript Tag 5}
  \end{minipage}\hfill
  \begin{minipage}{0.4\linewidth}
    \raggedleft
    {\sffamily\footnotesize\color{lpGray}Dozent:}\\
    {\sffamily\small\color{lpInk}Can Siebert}\\[10pt]
    {\sffamily\footnotesize\color{lpGray}Niveau-Mix:}\\
    {\sffamily\small\color{lpInk}3\,L1 \textperiodcentered{} 5\,L2 \textperiodcentered{} 3\,L3}
  \end{minipage}

  \vspace{0.5cm}
  \noindent{\color{lpAccent}\rule{\linewidth}{0.6pt}}
\end{titlepage}

\section*{Hinweise zu den Musterl\"osungen}
\thispagestyle{plain}

\noindent Die folgenden Musterl\"osungen sind \emph{eine} m\"ogliche, didaktisch nachvollziehbare L\"osung. Auf L2/L3 sind mehrere Begr\"undungswege legitim, solange sie:

\begin{itemize}
  \item den Begriffsapparat (Capability, Prozess, CIM/PIM/PSM, Modelling Charter) sauber nutzen,
  \item Annahmen explizit machen,
  \item Zielkonflikte (Standardisierung vs.\ Freiheit, Standard- vs.\ f\"oderiertes Tool) benennen,
  \item zu einer klar formulierten Entscheidung f\"uhren.
\end{itemize}

\begin{infobox}[Nutzung im Coaching]
Vergleichen Sie Ihre Antwort \emph{vor} der Musterl\"osung. Wo weicht Ihre Argumentation ab \textendash{} ist das eine andere Annahme oder ein Denkfehler?
\end{infobox}

\clearpage

\section*{Niveau L1 \textendash{} Wissen \& Reproduktion}
\addcontentsline{toc}{section}{L1 -- Wissen}

\aufgabenkopf{1}{Visio/Lucidchart als Enterprise-Werkzeug}{\nivLi}{8\,min}

\ytquery{Lucidchart Visio Enterprise Kollaboration Shape Library Vergleich}

\begin{loesungbox}
\textbf{Vier Enterprise-Capabilities:}

\begin{enumerate}
  \item \emph{Kollaboration} \textendash{} gleichzeitige Bearbeitung, Kommentare, Review-Tracking. \emph{Anti-Muster ohne:} ``E-Mail-Ping-Pong mit Anh\"angen'' \textendash{} niemand wei{\ss}, welche Version aktuell ist.
  \item \emph{Standards} \textendash{} Templates, Shapesets, Naming-Konventionen, Modellhierarchie, ``Definition of Done''. \emph{Anti-Muster:} ``Jeder modelliert anders'' \textendash{} Vergleichbarkeit verloren.
  \item \emph{Freigaben} \textendash{} Draft / Review / Approved, \"Anderungs\-antr\"age, Rollen. \emph{Anti-Muster:} ``\"Anderungen passieren irgendwie'' \textendash{} keine Audit-F\"ahigkeit.
  \item \emph{Single Source of Truth} \textendash{} Verlinkung zu Policies, Arbeitsanweisungen, KPIs, Risiken. \emph{Anti-Muster:} ``Modelle leben isoliert'' \textendash{} Compliance-Kette bricht.
\end{enumerate}

\smallskip
\emph{Merksatz:} Das Tool ist nicht das Differenzial. Die \emph{Disziplin}, mit der die vier Capabilities eingesetzt werden, ist es.
\end{loesungbox}

\clearpage

\aufgabenkopf{2}{Enterprise Architecture \textendash{} f\"unf Sichten}{\nivLi}{10\,min}

\ytquery{Enterprise Architecture Capability Map Sichten TOGAF ArchiMate Beispiel}

\begin{loesungbox}
\begin{tabularx}{\linewidth}{@{}p{3.6cm}|X|X|X@{}}
\toprule
\textbf{Sicht} & \textbf{Was zeigt sie?} & \textbf{Zielgruppe} & \textbf{L\"ost\ldots} \\
\midrule
Capability Map & ``Was k\"onnen wir?'' \textendash{} F\"ahigkeiten, unabh\"angig von Prozess/Tool. & Vorstand, Strategie. & Priorisierung: wo investieren? \\
\midrule
Prozesslandkarte & ``Wie arbeiten wir?'' \textendash{} E2E-Prozesse, hierarchisch. & CPO, Process Owner. & Steuerung: wo Engpass? \\
\midrule
Applikations\-landschaft & ``Womit arbeiten wir?'' \textendash{} Systeme + Owner. & CIO, IT. & Konsolidierung: wo Redundanz? \\
\midrule
Datenobjekt-\"Ubersicht & Welche Daten\-objekte (Kunde, Vertrag\ldots), System of Record? & CDO, Data Steward. & Datenqualit\"at + Compliance. \\
\midrule
Schnittstellen-/Integrations-Map & Wer redet mit wem (API, Event, Datei)? & EA, Integration Lead. & Architekturpfad: Risiken / Br\"uche. \\
\bottomrule
\end{tabularx}

\smallskip
\emph{Zusammenspiel:} Capability Map zeigt \emph{Sollst\"arke}, Prozesslandkarte zeigt \emph{Ablauf}, Applikations\-landschaft zeigt \emph{Realit\"at}. Drei Sichten = drei Hebel.
\end{loesungbox}

\clearpage

\aufgabenkopf{3}{MDA \textendash{} CIM / PIM / PSM}{\nivLi}{8\,min}

\ytquery{Model Driven Architecture MDA CIM PIM PSM Beispiel OMG}

\begin{loesungbox}
\begin{tabularx}{\linewidth}{@{}p{2.0cm}|X|X@{}}
\toprule
\textbf{Ebene} & \textbf{Geh\"ort hinein} & \textbf{Geh\"ort \emph{nicht} hinein} \\
\midrule
\textbf{CIM} (Business-Ebene) & Ziele, Begriffe (Glossar), Stakeholder, Prozesse, Datenobjekte (logisch). Owner: Fachbereich / CPO. & Technologien, APIs, Datenbanken, Plattform-Namen. \\
\midrule
\textbf{PIM} (logisch) & Services (logisch), Schnittstellen-Vertr\"age, Fehlerf\"alle, Datenstr\"ome. Owner: Enterprise Architect. & Konkrete Produkte / Anbieter / Deployment-Topologie. \\
\midrule
\textbf{PSM} (technisch) & Konkrete Plattform, APIs, Deployments, Secrets, Sicherheits-Implementierung. Owner: IT-Team / Architecture. & Business-Begriffe (sollten l\"angst in CIM stehen). \\
\bottomrule
\end{tabularx}

\smallskip
\textbf{Warum Verantwortungs-Trennung?} Wer \emph{auf CIM-Ebene} sagt ``Kunde in 24\,h aktiv'', sagt eine Business-Anforderung \textendash{} und der Fachbereich tr\"agt die Verantwortung. Wer \emph{auf PSM-Ebene} sagt ``\"uber unseren IAM-Provider in 30\,Sekunden zu provisionieren'', tr\"agt die technische Verantwortung. Vermischen f\"uhrt dazu, dass Fachbereich technische Designs vorgibt und Technik Business-Ziele neu interpretiert \textendash{} beides ist Fehler.
\end{loesungbox}

\clearpage

\section*{Niveau L2 \textendash{} Anwendung \& Transfer}
\addcontentsline{toc}{section}{L2 -- Anwendung}

\aufgabenkopf{4}{Modelling Charter in 45 Minuten}{\nivLii}{25\,min}

\ytquery{Modelling Charter Process Standards Naming Convention BPM Governance}

\begin{loesungbox}
\textbf{1. F\"unf nicht verhandelbare Standards:}
\begin{enumerate}
  \item \emph{Naming} (Verb + Objekt; Prozessname ``X-to-Y''). \emph{Sanktion:} Modell bleibt in Draft.
  \item \emph{Start / Ende explizit modelliert}. \emph{Sanktion:} Review-Abweisung.
  \item \emph{Rollen sichtbar} (Lanes / Swimlanes / Org-Annotation). \emph{Sanktion:} Approval-Block.
  \item \emph{Status / Version im Metadaten-Header}. \emph{Sanktion:} Modell nicht recherchierbar im Repository.
  \item \emph{Owner benannt (Rolle, nicht Person)}. \emph{Sanktion:} kein Lifecycle-\"Ubergang.
\end{enumerate}

\smallskip
\textbf{2. Drei Freiheitsgrade:}
\begin{itemize}
  \item \emph{Detailtiefe} \textendash{} je nach Modellzweck (Management-\"Ubersicht vs.\ Ausf\"uhrung).
  \item \emph{Notation f\"ur \"Ubersichten} \textendash{} EPC, Capability Map, Wertstromskizze nach Wahl.
  \item \emph{Lokale Varianten als Attribute} \textendash{} Standort darf Sonderpfade pflegen, sofern dokumentiert.
\end{itemize}

\smallskip
\textbf{3. Freigabeprozess:}\\
Draft (Modeler) \textrightarrow{} Review (BPM-Lead) \textrightarrow{} Approved (Process Owner). R\"uckkante: Review \textrightarrow{} Draft bei Mangel. Retired bei Ablauf G\"ultigkeit.

\smallskip
\textbf{4. Entscheidung unter Unsicherheit:}
\begin{itemize}
  \item \emph{In 2 Wochen:} f\"unf Standards verabschieden, Vorlagen-Set publizieren, ein Pilot-Bereich startet.
  \item \emph{In 3 Monaten:} Tool-Migration, RACI-Rollout, KPI-Set live.
\end{itemize}
Begr\"undung: Quick Wins zuerst, Tool-Migration ben\"otigt Vorbereitung; Akzeptanz-Risiko sinkt, wenn die Charter sichtbar wirkt, bevor Tools wechseln.
\end{loesungbox}

\clearpage

\aufgabenkopf{5}{Capability Map ``Digital Customer Onboarding''}{\nivLii}{22\,min}

\ytquery{Capability Map Customer Onboarding Enterprise Architecture Hebel}

\begin{loesungbox}
\textbf{Capability Map (10 Capabilities, in 3 Cluster):}

\smallskip
\emph{Cluster ``Kunden-Identit\"at \& Vertrag''}: Kunde identifizieren \textperiodcentered{} Identit\"at pr\"ufen \textperiodcentered{} Vertrag erstellen \textperiodcentered{} Vertrag pr\"ufen.\\
\emph{Cluster ``Zugang \& Konto''}: Konto anlegen \textperiodcentered{} Zugang provisionieren \textperiodcentered{} Rollen vergeben.\\
\emph{Cluster ``Compliance \& Aktivierung''}: KYC/Compliance-Check \textperiodcentered{} Aktivierung \textperiodcentered{} Hypercare starten.

\smallskip
\textbf{Top-3 mit h\"ochstem Hebel:}
\begin{itemize}
  \item \emph{Identit\"atspr\"ufung} \textendash{} dominiert sowohl DLZ als auch Compliance.
  \item \emph{Zugang provisionieren} \textendash{} dominiert die Durchlaufzeit (typischer Engpass).
  \item \emph{KYC/Compliance-Check} \textendash{} hohes Risiko bei Fehlern.
\end{itemize}

\smallskip
\textbf{E2E-Prozesslandkarte (vereinfacht):}\\
Lead $\to$ Vertrag $\to$ KYC $\to$ Provisionierung $\to$ Aktivierung $\to$ Hypercare $\to$ Bestandsf\"uhrung.

\smallskip
\textbf{Capability-Zuordnung:} Identit\"at pr\"ufen \& KYC bei \emph{Vertrag}; Provisionierung bei \emph{Provisionierung}; Rollen vergeben bei \emph{Aktivierung}.

\smallskip
\textbf{Drei kritische Datenobjekte:} \emph{Kunde}, \emph{Vertrag}, \emph{Identit\"atsnachweis}.

\smallskip
\textbf{Zwei Risiken:}
\begin{itemize}
  \item \emph{Datenqualit\"at:} Adress-/Personendaten unstet \textrightarrow{} R\"uckspr\"unge.
  \item \emph{Datenschutz:} KYC-Dokumente m\"ussen revisionssicher gespeichert sein \textendash{} DSGVO + L\"oschpflicht.
\end{itemize}
\end{loesungbox}

\clearpage

\aufgabenkopf{6}{MDA-Pfad f\"ur eine priorisierte Capability}{\nivLii}{25\,min}

\ytquery{MDA CIM PIM PSM Beispiel Identity Service Onboarding Capability}

\begin{loesungbox}
\textbf{Capability: \emph{Identit\"atspr\"ufung}.}

\smallskip
\textbf{1. CIM:}
\begin{itemize}
  \item \emph{Business-Ziel:} ``Kunde ist in 24\,h identifiziert, validiert und Konto aktivierbar.''
  \item \emph{Prozess-Scope:} Vertrags\-eingang \textrightarrow{} Identit\"atsnachweis \textrightarrow{} Konto\-anlage.
  \item \emph{Datenobjekt Kunde \textendash{} Pflichtfelder:} Name, Geburtsdatum, Anschrift, Identit\"atsnachweis-ID, Status.
\end{itemize}

\smallskip
\textbf{2. PIM:}
\begin{itemize}
  \item Service \emph{VerifyIdentity}: Input = Kundendaten + Nachweis; Output = Result (verified / failed) + Fehlerkategorie.
  \item Service \emph{CreateAccount}: Input = verifizierte Kundendaten; Output = Account-ID.
  \item Service \emph{AssignRoles}: Input = Account-ID + Tarif; Output = Rollen-Set.
\end{itemize}
\emph{Datenfluss:} Verify $\to$ Create $\to$ Assign; Fehlerf\"alle: \emph{fehlende Dokumente} (R\"uckfrage), \emph{Identit\"at nicht verifizierbar} (manueller Pfad).

\smallskip
\textbf{3. PSM:}
\begin{itemize}
  \item \emph{IAM-Plattform} (bestehend): provisioniert Konto.
  \item \emph{Identity-Provider}: externe Schnittstelle (z.\,B.\ Video-Ident-Anbieter).
  \item \emph{Workflow-Engine}: orchestriert die Sequenz.
  \item \emph{Audit-Logging Pflicht}: jeder Schritt protokolliert mit Zeitstempel + Ergebnis.
\end{itemize}

\smallskip
\textbf{Schlusssatz:} Das Business-Ziel ``24\,h Identifizierung'' ist auf CIM-Ebene \emph{nicht verhandelbar} \textendash{} es ist der \emph{Erfolgsbegriff}. Welcher konkrete Identit\"atsanbieter (Video-Ident, eID, Foto-Upload) eingesetzt wird, darf auf PSM-Ebene noch ge\"andert werden \textendash{} solange das CIM-Ziel erreicht bleibt.
\end{loesungbox}

\clearpage

\aufgabenkopf{7}{Naming \& Konvention}{\nivLii}{15\,min}

\ytquery{Naming Convention BPM Files Best Practice Versioning Process Model}

\begin{loesungbox}
\begin{tabularx}{\linewidth}{@{}cp{4.0cm}lX@{}}
\toprule
 & \textbf{Original} & \textbf{Bewertung} & \textbf{Korrektur + Begr\"undung} \\
\midrule
A & \texttt{O2C\_final\_v3\_neu.vsdx} & schlecht & \texttt{O2C\_v1.3\_Draft\_DE}. ``final'' und ``neu'' sind Antimuster; Status geh\"ort in Metadaten. \\
B & \texttt{Bestellung} & schlecht & \texttt{Bestellung-erfassen\_v1.0\_Approved}. Substantiv ohne Verb, ohne Version. \\
C & \texttt{Order-to-Cash} & gut & passt; ggf.\ Version + Status erg\"anzen. \\
D & \texttt{ProcRevAB\_Hoehl\_2025.vsdx} & schlecht & \texttt{P2P-Review-Audit\_v2.0\_2025\_Approved}. Pers\"onlicher Marker entfernen. \\
E & \texttt{Rechnung erstellen v1.0 - Draft - DE} & gut & passt; ggf.\ Trennzeichen vereinheitlichen. \\
F & \texttt{Onboarding (cleaner version - DO NOT DELETE)} & schlecht & \texttt{Onboarding\_v2.1\_Approved}. ``DO NOT DELETE'' ist ein Lifecycle-Signal, geh\"ort in Status. \\
G & \texttt{P2P Subprozess - QA - 2026-06-02} & mittel & \texttt{P2P-Subprozess-QA\_v1.0\_Approved\_2026-06-02}. Status fehlt. \\
H & \texttt{Untitled Diagram (37)} & schlecht & sofort umbenennen oder l\"oschen; ``Untitled'' kommt nicht in Approved. \\
\bottomrule
\end{tabularx}
\end{loesungbox}

\clearpage

\aufgabenkopf{8}{KPI-Set f\"ur Modellqualit\"at}{\nivLii}{18\,min}

\ytquery{KPI Modellqualitaet Process Repository Goodhart Compliance Steuerung}

\begin{loesungbox}
\textbf{Vier KPIs:}

\begin{tabularx}{\linewidth}{@{}p{3.6cm}llX@{}}
\toprule
\textbf{KPI} & \textbf{Typ} & \textbf{Frequenz} & \textbf{Owner + Formel} \\
\midrule
Lead Time Idee~\textrightarrow{}~Approved-Modell & Outcome (Speed) & monatlich & BPM-Lead; Median (Approval - Draft). \\
Anteil Approved-Modelle mit aktuellem Review & Outcome (Akt.) & quartalsweise & Process Owner; Approved-Modelle mit Review < 12\,Mon / Approved-Modelle gesamt. \\
Audit-Findings zu Prozess-Dokumentation & Risk/Compliance & j\"ahrlich & Audit-Lead; Anzahl Findings. \\
Quote nicht konformer Modelle (GoM-Check) & Risk/Compliance & monatlich & BPM-Lead; nicht-konform / gesamt. \\
\bottomrule
\end{tabularx}

\smallskip
\textbf{Manipulationsrisiken:}
\begin{itemize}
  \item \emph{Lead Time:} Modelle zu fr\"uh approved (oberfl\"achlicher Review). \emph{Schutz:} Stichproben-Audit auf Modell-Qualit\"at.
  \item \emph{Anteil aktueller Reviews:} Reviews pro forma (``Datum gesetzt, Inhalt nicht gepr\"uft''). \emph{Schutz:} Stichprobe + Review-Protokoll-Pflicht.
  \item \emph{Audit-Findings:} ``Nichts berichten'' (Findings nicht eskalieren). \emph{Schutz:} externe Audit-Stichprobe.
  \item \emph{Konformit\"ats-Quote:} Definition aufweichen. \emph{Schutz:} GoM-Definition zentral pflegen.
\end{itemize}

\smallskip
\textbf{Kopplung:} Lead Time \emph{darf nicht} verbessert werden, wenn die Konformit\"ats-Quote sinkt \textendash{} sonst ``Schnellschuss-Approvals''.

\smallskip
\textbf{Nicht eingef\"uhrt:} \emph{``Anzahl Modelle pro Modeler''} \textendash{} f\"ordert Quantit\"at \"uber Qualit\"at und ist klassisches Anti-Pattern; wird durch Lead Time + Konformit\"ats-Quote bereits abgedeckt.
\end{loesungbox}

\clearpage

\section*{Niveau L3 \textendash{} Management \& Entscheidungsarchitektur}
\addcontentsline{toc}{section}{L3 -- Management}

\aufgabenkopf{9}{Fallstudie \emph{RetailGroup Europe}}{\nivLiii}{45\,min}

\ytquery{Enterprise Model Governance Capability Map MDA Roadmap Retail Audit}

\begin{loesungbox}
\textbf{1. Modelling Charter:} 5 Standards (Naming, Owner, Status, Rollen, Start/Ende), 3 Freiheitsgrade (Detailtiefe, lokale Varianten, Visual-Style innerhalb Template), Freigabefluss Draft\,/\,Review\,/\,Approved (siehe Aufgabe 4).

\smallskip
\textbf{2. Capability Map (10\,--\,12):} Kunde verstehen \textperiodcentered{} Filialbetrieb \textperiodcentered{} Online-Shop \textperiodcentered{} Sortimentssteuerung \textperiodcentered{} Lieferung \& Retouren \textperiodcentered{} Zahlungsabwicklung \textperiodcentered{} Identit\"at \& Konto \textperiodcentered{} Compliance \textperiodcentered{} Personal \textperiodcentered{} IT-Plattformen.\\
\emph{Priorit\"aten:} \emph{Identit\"at \& Konto} (Onboarding) und \emph{Compliance} (Audit in 10\,W).

\smallskip
\textbf{3. MDA-Pfad f\"ur \emph{Identit\"at \& Konto}:} siehe Aufgabe 6 \textendash{} CIM ``Kunde in 24\,h aktiv'', PIM 3 Services, PSM via bestehendes IAM + Workflow-Engine.

\smallskip
\textbf{4. Vier KPIs:}
\begin{enumerate}
  \item \emph{Lead Time Onboarding} (Outcome / Speed): Median Start\,\textrightarrow{}\,Aktiv. Owner: COO.
  \item \emph{First-Pass-Yield Onboarding} (Outcome / Quality): \% F\"alle ohne R\"ucksprung. Owner: Onboarding-Lead.
  \item \emph{Compliance-Findings} (Risk): Anzahl pro Quartal. Owner: Compliance-Lead.
  \item \emph{Rework-Rate} (Risk / Quality): \% F\"alle mit Nacharbeit. Owner: QM-Lead.
\end{enumerate}

\smallskip
\textbf{5. Management\-entscheidung:}
\begin{itemize}
  \item \emph{Quick Wins Wochen 1\,--\,4:} (a) Modelling Charter publizieren (5 Standards), (b) Capability Map konsolidieren, mit Auditor abstimmen.
  \item \emph{Strukturell Wochen 5\,--\,10:} (c) MDA-Pfad ``Identit\"at'' bis PIM dokumentiert, (d) ein Pilot-Bereich auf das Tool migriert.
\end{itemize}

Begr\"undung: Quick Wins reduzieren \emph{Audit-Risiko sofort} (Charter + Capability Map = nachvollziehbare Struktur); strukturelle Ma{\ss}nahmen zeigen \emph{Steuerungs\-f\"ahigkeit} im Pilotbereich. Beide ohne neues Tool m\"oglich, passen in 100\,k-Budget.
\end{loesungbox}

\clearpage

\aufgabenkopf{10}{Tool-Strategie}{\nivLiii}{25\,min}

\ytquery{BPM Tool Strategy Standard vs Federated Visio Lucidchart Enterprise Decision}

\begin{loesungbox}
\textbf{1. Trade-off-Matrix:}

\begin{tabularx}{\linewidth}{@{}p{3.6cm}|X|X@{}}
\toprule
\textbf{Dimension} & \textbf{Standardtool} & \textbf{F\"oderiert} \\
\midrule
Kollaboration   & hoch (alle im selben Tool) & mittel (Inter-Tool-\"Ubersetzung n\"otig) \\
\midrule
Vergleichbarkeit & sehr hoch (1 Schema) & gering (Konvertierungsverlust) \\
\midrule
Akzeptanz       & mittel (Tool wird ``vorgeschrieben'') & hoch (Teams behalten Gewohnheiten) \\
\midrule
Migrations\-aufwand & hoch (alle Modelle umziehen) & gering (status quo plus Charter) \\
\bottomrule
\end{tabularx}

\smallskip
\textbf{2. Entscheidung:} \emph{Standardtool} \textendash{} aber mit Migrations-Welle.\\
Begr\"undung: Vergleichbarkeit + Audit-F\"ahigkeit \"uberwiegen den Akzeptanz-Schmerz. F\"oderiert spart kurzfristig Aufwand, kostet aber dauerhaft (Konvertierungsverluste, Inkonsistenzen, doppelte Pflege). Akzeptanz kann durch Pilot-Bereiche + Coaching gesch\"utzt werden.

\smallskip
\textbf{3. Bewusst hingenommenes Risiko:} \emph{Akzeptanzproblem in den ersten 3 Monaten} \textendash{} Power-User m\"ussen umlernen. \emph{Fr\"uhwarnsignal:} Wenn nach 6 Wochen weniger als 30\,\% der priorisierten Modelle migriert sind \textrightarrow{} Coach-Intensit\"at verdoppeln, Coaching-Mittel umschichten.

\smallskip
\textbf{4. Migrationspfad (drei unumgehbare Schritte):}
\begin{enumerate}
  \item \emph{Tool-Anforderungs-Check} (BPMN-Export-F\"ahigkeit, Rollenmodell, Audit-Log) bestanden.
  \item \emph{Pilot-Bereich migriert} (1 Bereich, 10 Kernmodelle, sauber).
  \item \emph{Trainings-Curriculum} (Level 1\,--\,3) etabliert; nur dann Welle 2.
\end{enumerate}
\end{loesungbox}

\clearpage

\aufgabenkopf{11}{Transferprojekt \textendash{} Enterprise Modelling Operating Model}{\nivLiii}{45\,min}

\ytquery{Enterprise Modelling Operating Model EA MDA Governance Rollout CEA}

\begin{loesungbox}
\textbf{1. Top-8 Entscheidungen:}
\begin{enumerate}
  \item Priorisierung der zu modellierenden Capabilities/Prozesse.
  \item Standardisierung (Naming, Charter, Templates).
  \item Investitionen (Tool, Schulungen, externer Coach).
  \item Risiken (Compliance, Audit-Findings).
  \item Schnittstellen (API-Standards, Integration).
  \item Tool-Strategie (Standard vs.\ f\"oderiert).
  \item Variantenstrategie (Attribute vs.\ Subprozesse).
  \item Reporting (KPI-Set, Reviews).
\end{enumerate}

\smallskip
\textbf{2. Modelling Charter v2:} 5 nicht verhandelbare Standards + 3 Freiheitsgrade + Lifecycle + KPI-Set (siehe Aufgabe 8). Erg\"anzung: \emph{Approval-Signatur ist Pflicht}; \emph{Modelle ohne Owner kommen nicht aus Draft}.

\smallskip
\textbf{3. EA-Minimum Viable Architecture:}
\begin{itemize}
  \item Capability Map (10\,--\,12) konsolidiert.
  \item Prozesslandkarte (5 Kernprozesse) verkn\"upft mit Capabilities.
  \item Applikations-/Datenobjekt-Minimum: 5 zentrale Applikationen, 5 zentrale Datenobjekte, jeweils mit \emph{Data Owner}.
\end{itemize}

\smallskip
\textbf{4. MDA-Guardrails:}
\begin{itemize}
  \item CIM ist \emph{verbindlich} f\"ur Ziele, Begriffe, Datenpflichtfelder.
  \item PIM darf Services / Schnittstellen / Fehlerf\"alle definieren, aber \emph{keine} Plattform-Wahl treffen.
  \item PSM ist iterativ \"anderbar, solange CIM-Ziele erreicht bleiben.
  \item Review-Gate: jede neue Capability durchl\"auft CIM-Review (Fachbereich) bevor PIM beginnt.
\end{itemize}

\smallskip
\textbf{5. Rollout \& Enablement:}
\begin{itemize}
  \item Welle 1 (Pilot, Wo 1\,--\,4): 1 Bereich, 10 Modelle, ein MDA-Pfad.
  \item Welle 2 (Scale, Wo 5\,--\,8): 3 Bereiche, Tool-Migration startet.
  \item Welle 3 (Sustain, Wo 9\,--\,12): KPI-Reviews monatlich, Governance Board konstituiert.
  \item Trainings: L1 (Lesen), L2 (Modellieren), L3 (Architektur \& Entscheidungen).
  \item Akzeptanzstrategie: Power-User als Multiplikatoren; Pilot-Erfolg sichtbar machen.
\end{itemize}

\smallskip
\textbf{6. Zwei Entscheidungen unter Unsicherheit:}
\begin{enumerate}
  \item \emph{Tool-Strategie:} Standardtool mit Migrations-Welle (siehe Aufgabe 10). \emph{Risiko:} Akzeptanz; \emph{Fr\"uhwarnsignal:} $<$\,30\,\% Migrationsfortschritt nach 6\,W.
  \item \emph{Governance-Strenge:} \emph{Minimal-mit-Z\"ahnen} \textendash{} 5 nicht verhandelbare Standards, aber sonst pragmatisch. \emph{Akzeptanzplan:} Charter zusammen mit Power-Usern formulieren. \emph{Trigger zum Nachsch\"arfen:} $\geq 3$ Audit-Findings zu Prozess-Dokumentation in 6\,Monaten.
\end{enumerate}
\end{loesungbox}

\clearpage

\section*{Abschluss}
\thispagestyle{plain}

\noindent Die Musterl\"osungen sind eine Diskussionsbasis. Pr\"ufen Sie Ihre eigene L\"osung kritisch: Habe ich Capability und Prozess sauber unterschieden? Habe ich Entscheidungs-Verantwortung zwischen Fachbereich (CIM) und IT (PSM) klar getrennt?

\medskip
\begin{infobox}[N\"achster Schritt]
Bringen Sie Ihre L\"osung ins Coaching mit: Welche zwei Standards sind in Ihrer Realit\"at \emph{unverhandelbar}? Welche Capability ist Ihr Hebel? Welche MDA-Trennung haben Sie real durchgesetzt?
\end{infobox}

\vspace{1cm}
\noindent{\color{lpAccent}\rule{\linewidth}{0.6pt}}\\[4pt]
{\footnotesize\sffamily\color{lpGray}Lernplatt \textperiodcentered{} by Can Siebert \textperiodcentered{} Kurs 71 (Dig 42) \textperiodcentered{} Tag 5 \textperiodcentered{} L\"osungsheft \textperiodcentered{} \textcopyright{} 2026}

\end{document}
